% sesmesh — Solvent Excluded Surface meshes for molecules
% Copyright (c) 2026 Ugo Varetto
% SPDX-License-Identifier: MIT OR Apache-2.0

\section{Availability}
\label{sec:availability}

\begin{center}
\small
\begin{tabular}{@{}>{\bfseries}l p{0.7\linewidth}@{}}
  Software & \sesmesh{} 1.1.2, released 2026-11-15 (first public release
    0.9.0 and 1.0.0 on 2026-11-13, 1.1.0 and 1.1.1 on 2026-11-14); library crate \code{sesmesh}, command-line
    crate \code{sesmesh-cli} (binary \code{sesmesh})~\cite{sesmesh} \\
  Source code & \url{https://github.com/uvar-apps/sesmesh} \\
  Packages & \url{https://crates.io/crates/sesmesh},
    \url{https://crates.io/crates/sesmesh-cli};
    API documentation at \url{https://docs.rs/sesmesh} \\
  Binaries & GitHub releases for macOS arm64, Linux x86-64 and arm64,
    Windows x86-64 \\
  Website & \url{https://uvar-apps.github.io/sesmesh/} (guides, tutorials,
    gallery, interactive viewer) \\
  Archive & Zenodo, version 1.1.2: \href{https://doi.org/10.5281/zenodo.24238837}{doi:10.5281/zenodo.24238837}\newline
    all versions: \href{https://doi.org/10.5281/zenodo.24238836}{doi:10.5281/zenodo.24238836} \\
  Licence & MIT or Apache-2.0, at the user's option \\
  Language & Rust 1.89 or newer; no unsafe code; dependencies from
    crates.io only \\
  Install & \code{cargo install sesmesh-cli}; library: \code{sesmesh = "1"}
    in \code{Cargo.toml} \\
\end{tabular}
\end{center}

\paragraph{Data and reproducibility.}
The validation data used in this paper are in the repository
(\code{docs/validation/}: accuracy study, robustness sweeps, large inputs,
fuzzing campaign),
together with the commands that regenerate them. Structures are fetched from
the RCSB PDB by \code{cargo xtask fetch}, whose lock file pins their SHA-256
checksums; PDB data are in the public domain (CC0). The paper's sources,
figures, the script that turns the data into its tables and plots, and the
program that produces the analytic-reference data are in \code{paper/}
(\cref{app:reproduce}).

\section*{Declarations}

\paragraph{Use of AI tools.}
The software, its tests and documentation, and a first draft of this manuscript,
except for some of the benchmarks, were developed with the help of an AI coding assistant 
(Anthropic's Claude Opus 5.5, used through Claude Code).
The assistant also ran the validation pipeline locally and, at the author's
request, made the commits, release tags and package publications under the
author's own Git identity; every commit to which it contributed carries a
\code{Co-Authored-By: Claude} trailer in the project's version history.
The author designed and directed the work, implemented draft versions of some
of the algorithms in Python and C++, reviewed the code, the measurements and 
the text, and takes full responsibility for the content.

\paragraph{Affiliations.}
The author is also affiliated with CSIRO, the Pawsey Supercomputing Research
Centre and The University of Western Australia. This work was carried out
independently of those roles, on the author's own time and equipment and at
the author's own expense; it received no support from these institutions and
does not represent their views.

\paragraph{Competing interests.}
The author declares no competing interests.

\paragraph{Funding.}
This research did not receive any specific grant from funding agencies in
the public, commercial, or not-for-profit sectors; the author bore all costs.

\paragraph{Author contributions.}
Ugo Varetto is the sole author: conceptualization, methodology, software,
validation, formal analysis, investigation, data curation, visualization,
writing (original draft, review and editing).

\paragraph{Declaration of generative AI and AI-assisted technologies in the manuscript preparation process.}
During the preparation of this work the author used Claude Opus 5.5
(Anthropic), through Claude Code, to write a first draft of the manuscript.
After using this tool, the author reviewed and edited the content as needed
and takes full responsibility for the content of the published article.

\paragraph{Acknowledgements.}
The atomic coordinates used in this work were obtained from the 
Protein Data Bank~\cite{berman2000pdb} through the
RCSB PDB (\url{https://www.rcsb.org}), a member of the worldwide PDB
(wwPDB, \url{https://www.wwpdb.org})~\cite{berman2003wwpdb,wwpdb2019}.
Entries are identified by their PDB IDs; the full lists are in
\code{tests/data/sets.toml}.
